% Recovery of the native product and its selector, not incorporation of RH.
\subsection{Multiplicative irreducibility and return periods}
\label{primos-retornos}

The distinction among structure, evaluation, and phase also appears in the
multiplicative sector. The local APP product contains both the remainder
and the quotient; its iteration makes it possible to construct normal forms
and factorizations before examining their decimal periods. This development
places the expression \emph{prime mode} in a precise domain and relates
irreducibility to the operations already defined, not to an isolated
coincidence of periods.

Let
\[
 \mathcal W_9=\{\varnothing\}\sqcup
 \bigsqcup_{\ell\geq1}\{1,\ldots,9\}^{\ell}.
\]
The sequences are ordered from the least significant digit. The evaluation
\(\operatorname{ev}_9(u)=\sum_j u_j9^j\), with
\(\operatorname{ev}_9(\varnothing)=0\), is bijective onto the nonnegative
integers: for a positive integer, division of \(n-1\) by \(9\)
uniquely determines the first digit and the continuation. This is the
bijective nonadic numeration; it preserves the representative \(9\) of the
boundary.

The sum \(\boxplus_\#\) is obtained by applying the additive sheet of APP to
the digits at each position and propagating their quotients. If there is only
one digit, it constitutes the provisional remainder; a nonzero incoming
quotient is normalized by a second additive operation. The sum of the
quotients is transported to the next position. To multiply a sequence by a
digit \(d\), the multiplicative cell produces \(u_jd=9q_j+r_j\), and the
incoming quotient \(c_j\) is incorporated through
\[
 (w_j,c_{j+1})=
 \begin{cases}
 (r_j,q_j),&c_j=0,\\
 \bigl(R^+(r_j,c_j),q_j+q^+(r_j,c_j)\bigr),&c_j>0.
 \end{cases}
\]
One begins with \(c_0=0\), and the final quotient is incorporated when it is
nonzero. During addition, \(c_j\leq2\); during multiplication by a digit,
\(c_j\leq9\). The local operations thus receive arguments within
their domain. Denote this second procedure by \(d\odot_\#u\),
with \(d\odot_\#\varnothing=\varnothing\).

For \(v=(v_0,\ldots,v_{m-1})\), the full product is constructed by
the recurrence
\[
 a_m=\varnothing,\qquad
 a_j=(9\odot_\#a_{j+1})\boxplus_\#(v_j\odot_\#u),
 \qquad u\boxtimes_\#v=a_0.
\]
The rule composes the two local sheets. Its evaluation verifies
\begin{equation}
 \operatorname{ev}_9(u\boxtimes_\#v)
   =\operatorname{ev}_9(u)\operatorname{ev}_9(v).
 \label{eq:primos-entrelazamiento}
\end{equation}
Indeed, each transition preserves the partial sum plus the pending quotient
weighted by the corresponding power of \(9\). The recurrence then gives
\(\operatorname{ev}_9(a_j)=9\operatorname{ev}_9(a_{j+1})+
v_j\operatorname{ev}_9(u)\), whose iteration proves
\eqref{eq:primos-entrelazamiento}. The uniqueness of the normal form transports
associativity and commutativity to the constructed product; its unit is \((1)\).

\begin{definition}[Reduced multiplicative coproduct]
On the finite-support vector space generated by the nonempty normal
forms, define
\[
 \widetilde\Delta_\# e_w
 =\sum_{\substack{u\boxtimes_\#v=w\\u,v\ne(1)}}e_u\otimes e_v.
\]
A nonunit normal form is irreducible when its image under
\(\widetilde\Delta_\#\) is zero.
\end{definition}

The factorization fibers are finite by
\eqref{eq:primos-entrelazamiento}. In the completion
\(\mathcal H_\#=\ell^2(\mathcal W_9\setminus\{\varnothing\})\),
the images of distinct forms have disjoint supports. If
\(d_\#(w)\) counts the ordered nonunit factorizations, the domain
of the closure is
\[
 \operatorname{Dom}(\overline{\widetilde\Delta_\#})
 =\left\{x\in\mathcal H_\#:
       \sum_w d_\#(w)|x_w|^2<\infty\right\}.
\]
\Needspace{8\baselineskip}
Indeed, the squared norm of the image is that weighted sum. Therefore,
the operator
\[
 A_\#=(\overline{\widetilde\Delta_\#})^*
                  \overline{\widetilde\Delta_\#}
\]
is diagonal, positive, and self-adjoint: its eigenvalue at \(e_w\) is the
number of ordered nonunit factorizations of \(w\). The projection
\[
 P_{\rm irr}=(I-|e_{(1)}\rangle\langle e_{(1)}|)
                    \mathbf1_{\{0\}}(A_\#)
\]
selects exactly the irreducible forms. Through the multiplicative
evaluation, these correspond to the ordinary primes. The selection
depends on the factorization structure and keeps the unit state separate.

For a value \(n\geq2\) coprime to \(10\), decimal division induces the
permutation of remainders \(r\mapsto10r\pmod n\). The return of remainder
\(1\) has period
\[
 \tau_n=\operatorname{ord}_n(10).
\]
If \(n=\prod_p p^{a_p}\), the Chinese remainder theorem gives
\[
 \tau_n=\operatorname{lcm}_{p^{a_p}\parallel n}
                    \operatorname{ord}_{p^{a_p}}(10).
\]
The composite return synchronizes the returns of the prime powers.
For a prime \(p\notin\{2,3,5\}\), one has
\[
 \tau_p\mid p-1,\qquad
 M_p=\frac{10^{\tau_p}-1}{p}\in\mathbb N,
 \qquad M_p\equiv0\pmod9.
\]
The last congruence expresses the nonadic closure of the decimal repetition:
\(10^{\tau_p}-1\) is a multiple of \(9\), and \(p\) is invertible modulo
\(9\). The same congruence holds for every \(n\) coprime to
\(90\). In particular, \(7,13,77,91,143\) have decimal period \(6\),
although only the first two are irreducible. The coproduct and the period
therefore record different properties of the same constructed value.

This differentiation specifies the relation with aperiodic refinement.
The vacancy sequence determines the increments of the capacity index;
factorization determines the multiplicative modes; decimal division
determines their return times. The nonadic representative preserves phase
closure, whereas the quotient and factorization preserve information that
this closure does not distinguish. The link with the spectral constructions
of the treatise resides in those operators and their intertwining maps. The
proof of its subsequent analytic results retains its own development; it is
not replaced by the periodicity of a residual observation.

\subsubsection{Vacancy discrepancy and Dirichlet readings}

The link also admits an exact analytic expression. The same coding
\(z_n=z_n(0)\), for \(n\geq1\), is retained, and two subsequent
observables are considered: the series over all indices and the product
over the indices selected by irreducibility. For
\(\operatorname{Re}s>1\),
\[
 Z_{\rm vac}^{+}(s)=\sum_{n\geq1}\frac{z_n}{n^s},\qquad
 Z_{\rm vac}^{\times}(s)=\prod_p(1-p^{-s})^{-z_p}.
\]
Both converge absolutely in that half-plane. Their common factor is the
density \(\delta\), but their discrepancies are different.

\Needspace{21\baselineskip}
\begin{proposition}[Separation of the vacancy density]
One has
\[
 Z_{\rm vac}^{+}(s)=\delta\zeta(s)+H_{\rm vac}(s),
 \qquad H_{\rm vac}\in\mathcal O(\operatorname{Re}s>0).
\]
In \(\operatorname{Re}s>1\), with logarithms defined by the absolutely
convergent Euler series,
\[
 Z_{\rm vac}^{\times}(s)=\zeta(s)^\delta G_{\rm vac}(s),\qquad
 \log G_{\rm vac}(s)=
 \sum_p(z_p-\delta)\log(1-p^{-s})^{-1}.
\]
\end{proposition}
\begin{proof}
The centered partial sums satisfy
\[
 A_N:=\sum_{n=1}^N(z_n-\delta)
 =\lfloor(N+1)\delta\rfloor-N\delta,\qquad |A_N|<1.
\]
Abel summation represents the centered term by
\(s\int_1^\infty A_{\lfloor x\rfloor}x^{-s-1}\,dx\). The integral
converges locally uniformly in \(\operatorname{Re}s>0\),
which proves holomorphy. The second identity results from separating
\(z_p=\delta+(z_p-\delta)\) in the logarithm of the product, within
its domain of absolute convergence.
\end{proof}

The decomposition makes visible what the reading over primes adds: it
examines the same vacancy sequence on a subset determined by
factorization. Controlling the discrepancy over all integers gives
the first holomorphic continuation; the second observable preserves a
discrepancy of its own over primes. These are two applications of the same
discrete structure, with explicit analytic domains. This articulation
establishes its relevance to the arithmetic core without transferring to
this article the proof of the terminal results of the treatise.
